\section{Compact Hamiltonian lattice QED with exact Gauss law}
\label{lat-section}

We begin with a local ultraviolet regulator. This section proves a vacuum-relative
bound for a spatially local, gauge-invariant Hamiltonian with compact
$U(1)$ link variables and Wilson fermions.  The electric rotors are not
truncated.  The statement is about this specified Hamiltonian, with a
nonnegative bare Wilson mass; it is not a continuum-limit assertion.
The Hamiltonian formulation follows the framework of
Kogut--Susskind~\cite{KogutSusskind}.
For the direct spatial discretization of the Dirac Hamiltonian with a
Wilson term, see also
Zache et al.~\cite{Zache}, Eq.~(3).
The proof below establishes the comparison directly, without importing a
Euclidean transfer-matrix assertion from a differently discretized action.

\subsection{Hamiltonian, signs, and physical sectors}

Let $X=(\mathbb Z/n\mathbb Z)^3$, $n\ge2$, have spacing $a>0$ and
$N=n^3$ sites.  There is a positively oriented link $(x,j)$ from $x$ to
$x+e_j$ for every $x$ and $j=1,2,3$.  Write $\mathcal E$ for these
$3N$ links and set
\begin{equation}
 \mathcal H=L^2(U(1)^{\mathcal E},dU)
       \otimes\mathcal F_{\!a}(\ell^2(X)\otimes\mathbb C^4).
 \label{lat-hilbert}
\end{equation}
All Haar measures are normalized.  On a link,
$U_{xj}=e^{i\varphi_{xj}}$ acts by multiplication and
$E_{xj}=-i\partial_{\varphi_{xj}}$ has spectrum $\mathbb Z$.
For $g^2>0$, let
\begin{align}
 H_E&=\frac{g^2}{2a}\sum_{x,j}E_{xj}^2,\qquad
 H_B(U)=\frac1{g^2a}\sum_p(1-\operatorname{Re}U_p),
 \label{lat-gauge-energy}\\
 U_p&=U_{xj}U_{x+e_j,k}\overline{U_{x+e_k,j}}\,
          \overline{U_{xk}},\qquad j<k.
 \nonumber
\end{align}
Each elementary plaquette occurs once.  No electric background shift,
topological angle, or extra global-flux projection is imposed.

Fix four Hermitian Dirac matrices $\alpha_1,\alpha_2,\alpha_3,\beta$,
which square to one and anticommute pairwise.  The convention for a
negatively charged fermion is
\begin{align}
 (T_j(U)f)(x)&=\overline{U_{xj}}f(x+e_j),\qquad
 v_j(U)=\frac{T_j(U)-T_j(U)^*}{2ia},\label{lat-hopping}\\
 M(U)&=m+\frac r{2a}\sum_{j=1}^3
             \bigl(2-T_j(U)-T_j(U)^*\bigr),\qquad m\ge0,\quad r\ge0,
 \label{lat-wilson-mass}\\
 h_W(U)&=\sum_j\alpha_j v_j(U)+\beta M(U),\qquad
 H=H_E+H_B(U)+d\Gamma(h_W(U)).\label{lat-hamiltonian}
\end{align}
Spatial operators commute with the Dirac matrices.  Every $T_j(U)$ is
unitary, and consequently
\begin{equation}
 M(U)\ge mI,\qquad
 \|h_W(U)\|\le h_*:=m+\frac{3+6r}{a}.
 \label{lat-mass-positive}
\end{equation}
Indeed, $2-T-T^*=(I-T)^*(I-T)$.  This is the sign of the Wilson
term needed in the theorem.  The choice $r>0$ is the usual Wilson choice;
$r=0$ also obeys the proof but retains naive fermion doublers.

Put $n_x=\sum_{s=1}^4 c_{xs}^*c_{xs}$ and $q_x=n_x-2$.
Thus $q_x\in\{-2,-1,0,1,2\}$ measures electron excess, while physical
fermionic charge in units of $g$ is $-q_x$.  The centered charge is a
definition, not a subtraction chosen from the interacting state.
For $g_x=e^{i\theta_x}$, distinguish the gauge function $g_x$ from the
coupling $g$ and define
\begin{align}
 U^g_{xj}&=g_xU_{xj}\overline{g_{x+e_j}},\qquad
 (R_g\Psi)(U)=\Psi(U^g),\nonumber\\
 W(g)&=\exp\Bigl(i\sum_x\theta_xq_x\Bigr)R_g,
 \qquad
 G_x=\sum_j(E_{xj}-E_{x-e_j,j})+q_x.
 \label{lat-gauge-unitary}
\end{align}
The symbol $g$ in $R_g,W(g),U^g$ denotes the gauge function.
Because $h_W(U^g)=g^*h_W(U)g$, conjugation by the fermion factor
$\Gamma(g)$ cancels conjugation by $R_g$.  Hence $W(g)$ commutes
with all three terms of $H$.  Also
$W(g)=\exp(i\sum_x\theta_xG_x)$.
The $G_x$ strongly commute and have integer joint spectrum.

An external source is an integer function $\nu:X\to\mathbb Z$.
Positive nuclei have
\begin{equation}
 \nu_x=\sum_{i=1}^K Z_i\mathbf1_{\{x=R_i\}},\qquad
 Z_i\in\{1,\ldots,Z_{\max}\}.
 \label{lat-sources}
\end{equation}
The physical Hilbert space and its orthogonal projection are
\begin{align}
 \mathcal H_\nu&=\{\Psi:G_x\Psi=\nu_x\Psi\text{ for all }x\},\nonumber\\
 P_\nu&=\int_{[0,2\pi)^N}\frac{d\theta}{(2\pi)^N}
            e^{-i\sum_x\theta_x\nu_x}W(e^{i\theta}).
 \label{lat-projection}
\end{align}
Gauss's law is therefore
$\operatorname{div}E=\nu-q$: positive nuclei are screened by electron
excess.  In particular, on the torus
\begin{equation}
 \sum_xq_x=\sum_x\nu_x,\qquad
 \mathcal H_\nu\ne\{0\}\quad\Longleftrightarrow\quad
             -2N\le\sum_x\nu_x\le2N.
 \label{lat-capacity}
\end{equation}
For sufficiency, select site occupations with the required total excess.
The integer function $\nu-q$ has zero sum.  On a connected finite graph
it is the divergence of an integer link flow: route each site's integer
excess to a fixed root along a spanning tree.  The corresponding electric
Fourier vector and occupation vector give a state in $\mathcal H_\nu$.
Necessity follows by summing Gauss's law.  Integrality of $\nu$ is
essential for a character of compact $U(1)$; fractional external charges
would require a different representation or boundary prescription.

The finite fermion space still contains pair creation.  For $m>0$, the
free matrix $h_W(1)$ has $2N$ positive and $2N$ negative eigenvalues.
The finite particle--hole transform relative to its negative projection
identifies $\sum q_x$ with $N_e-N_p$.  The blocks
$P_+h_W(U)P_-$ of the interacting bilinear produce electron--positron
creation and annihilation.  No projection removes them.  The compact
gauge-field Hilbert space is infinite dimensional; it is not being
replaced by a classical link configuration in the Hamiltonian.
The Hamiltonian uses the raw bilinear \(d\Gamma(h_W(U))\).
In a particle--hole representation relative to \(P_-(h_W(1))\), its
generally link-dependent scalar
\(\Tr(P_-(h_W(1))h_W(U))\) is retained. Subtracting that function
of the links would define a different operator; it is not a harmless
constant vacuum shift. A single source-independent scalar constant
may be changed without changing the relative theorem.

\begin{lemma}[Operator realization and extensive lower bound]
\label{lat-realization}
The Hamiltonian \eqref{lat-hamiltonian} is self-adjoint on $D(H_E)$,
has compact resolvent, and $e^{-tH}$ is trace class for every $t>0$.
Every admissible physical restriction has these properties.  Moreover,
\begin{equation}
 H\ge-2Nh_*,\qquad
 -2Nh_*\le E_0\le-2N(m+3r/a)+\frac{3N}{g^2a},
 \label{lat-extensive}
\end{equation}
where $E_\nu=\inf\sigma(H|_{\mathcal H_\nu})$.
\end{lemma}
\begin{proof}
$H_B+d\Gamma(h_W)$ is a bounded continuous Hermitian matrix multiplier
on a finite product of circles.  Bounded perturbation of the positive
torus Laplacian proves self-adjointness and compact resolvent.
Eigenvalue comparison with $H_E$ proves trace-class heat semigroups.
The matrix $C_5=\alpha_1\alpha_2\alpha_3\beta$ anticommutes with
$h_W(U)$, so its eigenvalues occur in opposite pairs.  The smallest
eigenvalue of $d\Gamma(h_W(U))$ is
$-\tfrac12\operatorname{Tr}|h_W(U)|\ge-2Nh_*$.
The physical projections reduce $H$.
For the upper bound, take the constant link wave function and occupy
the two $\beta=-1$ spinors at every site.  This state has $q_x=0$,
zero electric field, and lies in $\mathcal H_0$.  Hopping expectations
vanish, its fermionic energy is $-2N(m+3r/a)$, and its magnetic energy
is at most $3N/(g^2a)$.  This gives the stated trial bound.
\end{proof}

\subsection{A positive trace for time-dependent Wilson matrices}

\begin{lemma}[Local twists and varying positive mass]
\label{lat-determinant}
Let $h(t)=\sum_j\alpha_j v_j(t)+\beta M(t)$ be bounded and measurable
on a finite time interval of length $b>0$, with Hermitian spatial
coefficients, $M(t)\ge mI$ and $m>0$.  Let $V$ be its ordered
one-particle propagator and $g=e^{i\Theta}$ any spin-scalar unitary,
where $\Theta$ is a Hermitian spatial matrix.  Then
\begin{equation}
 z(g,V):=e^{-\frac i2\operatorname{Tr}\Theta}\det(I+gV)>0.
 \label{lat-positive-determinant}
\end{equation}
Here the trace in the phase includes the four spin components.  The
inequality is nonnegative when $M(t)\ge0$ is obtained as a limit of
strictly positive masses.
\end{lemma}
\begin{proof}
Represent the propagator and the terminal unitary as the solution of
$U'=-A(t)U$, with $A=c(t)h(t)+i\Phi(t)$, $c\ge0$, positive total
free time, and Hermitian spin-scalar $\Phi$.  Free and kick intervals
can be separate; the kick $e^{i\Theta}$ has $\int\Phi=-\Theta$.
Since $\operatorname{Tr}h=0$,
$z=e^{\frac12\int\operatorname{Tr}A}\det(I+U)$ is the expression
in \eqref{lat-positive-determinant}.
The identities
\[
 A^*C_5+C_5A=0,\qquad U^*C_5U=C_5
\]
give
$\overline{\det(I+U)}=\det(U)^{-1}\det(I+U)$.
Together with $\det U=e^{-\int\operatorname{Tr}A}$, this makes $z$
real, including its centering phase.

If $Uv=-v$ for a nonzero vector, the associated solution $f$ is
antiperiodic and satisfies $f'+Af=0$.  But
\begin{equation}
 0=\operatorname{Re}\int f^*\beta(f'+Af)
      =\int c(t)f(t)^*M(t)f(t)\,dt>0.
 \label{lat-no-zero}
\end{equation}
The derivative contributes the zero antiperiodic boundary term;
$\beta\alpha_jv_j$ and $i\beta\Phi$ are anti-Hermitian.  The last
inequality holds because a nonzero solution never vanishes and the free
time has positive length.  Thus $z$ cannot be zero.
Continuously turn off $\Phi$ and $v_j$, and replace $M(t)$ by
$(1-s)M(t)+smI$.  This preserves the argument.  The terminal expression
is $\det(I+e^{-bm\beta})>0$, proving the sign.  Bounded measurable
coefficients follow either from the same absolutely continuous ODE
argument or from bounded step approximations.  Finally replace a
nonnegative $M(t)$ by $M(t)+\epsilon I$ and pass to the limit.
\end{proof}

For the local gauge twist, $\Theta_{xy}=\theta_x\delta_{xy}I_4$.
Therefore its phase is $e^{-2i\sum_x\theta_x}$, exactly the phase
from $q_x=n_x-2$.  It is single-valued under $\theta_x\mapsto
\theta_x+2\pi$.  Omitting this phase would invalidate the argument.

\subsection{Gauss-law domination with quantized compact links}

\begin{theorem}[Physical-sector vacuum comparison]
\label{lat-main}
For \eqref{lat-hamiltonian}, at any $a>0$, $g^2>0$, $r\ge0$, and
$m\ge0$, every admissible integer source $\nu$ satisfies, for all $b>0$,
\begin{equation}
 0<Z_\nu(b):=\operatorname{Tr}(P_\nu e^{-bH})
       \le Z_0(b),\qquad E_\nu\ge E_0.
 \label{lat-vacuum-domination}
\end{equation}
The zero-source physical vacuum is also a ground-energy sector of the
whole auxiliary Hilbert space \eqref{lat-hilbert}.  The assertion includes
the full electric energy of the sources; no nuclear self-energy has yet
been subtracted.
\end{theorem}
\begin{proof}
Set $c=g^2/(2a)$.  Relative to normalized Haar measure the one-link
electric heat kernel is
\begin{equation}
 k_t(\phi)=\sum_{\ell\in\mathbb Z}e^{-ct\ell^2}e^{i\ell\phi}
  =\sqrt{\frac\pi{ct}}\sum_{s\in\mathbb Z}
           e^{-(\phi-2\pi s)^2/(4ct)}>0.
 \label{lat-positive-heat}
\end{equation}
Let $K_t(U,V)$ be the product over links.  It is positive, continuous,
and obeys the Markov convolution identity.

For completeness, the bounded matrix Feynman--Kac formula used here
can be obtained directly.  Write $B(U)=H_B(U)+d\Gamma(h_W(U))$.
In the Dyson expansion for $e^{-b(H_E+B)}$, a $k$th-order kernel is
the time-simplex integral of $k$ matrices $B$ inserted between
$k+1$ electric heat kernels.  Its norm is at most
\begin{equation}
 K_b(U,V)\,\frac{(b\|B\|)^k}{k!}.
 \label{lat-dyson-kernel-bound}
\end{equation}
This follows by integrating the positive scalar kernels and using their
semigroup property.  The absolutely convergent series is the bounded
perturbation Dyson series and hence is the kernel of the required
semigroup.  The Markov bridge moments give the equivalent formula
\begin{equation}
 [e^{-bH}](U,V)=K_b(U,V)\,\mathbb E_{V\to U}^{b}
 \left[e^{-\int_0^bH_B(\omega_t)dt}
       \Gamma(V_\omega)\right].
 \label{lat-matrix-fk}
\end{equation}
$V_\omega$ is the ordered one-particle evolution with generator
$-h_W(\omega_t)$.  The bridge is the product of circle Brownian bridges.
The expression also follows by approximating its continuous paths with
step paths; boundedness of $B$ supplies domination.  Thus this use of
Feynman--Kac requires neither an unbounded matrix potential nor a
formal interchange of a nonconvergent expansion.

Inserting $W(g)$ changes the first kernel argument to $U^g$.
The trace therefore is
\begin{align}
 \operatorname{Tr}(W(g)e^{-bH})
   =\int dU\,K_b(U^g,U)\,\mathbb E_{U\to U^g}^{b}
     \left[e^{-\int_0^bH_B(\omega_t)dt}
        e^{-2i\sum_x\theta_x}\det(I+gV_\omega)\right].
 \label{lat-twisted-trace}
\end{align}
The finite-CAR identity $\operatorname{Tr}_{\mathcal F}\Gamma(A)
=\det(I+A)$ proves the determinant formula.  Traces of the continuous
kernels can be justified by factoring the semigroup into two
Hilbert--Schmidt factors; the resulting convolution gives the diagonal
kernel integral.  The bound \eqref{lat-dyson-kernel-bound}, with the
finite fermion trace, also bounds all integrals uniformly in the gauge
twist at fixed $b,n,a,g$.

For $m>0$ every integrand in \eqref{lat-twisted-trace} is strictly
positive: the scalar factors are positive and
Lemma~\ref{lat-determinant} applies to the time-dependent Wilson
matrix.  Denote this nonnegative function of $\theta$ by $F_b(\theta)$.
Equation~\eqref{lat-projection} gives
\begin{equation}
 Z_\nu(b)=\int\frac{d\theta}{(2\pi)^N}
          e^{-i\theta\cdot\nu}F_b(\theta),\qquad
 |Z_\nu(b)|\le\int\frac{d\theta}{(2\pi)^N}F_b(\theta)=Z_0(b).
 \label{lat-character-bound}
\end{equation}
For an admissible sector its actual heat trace is strictly positive,
so the absolute value may be removed.  Its compact-resolvent spectral
expansion gives
$-\lim_{b\to\infty}b^{-1}\log Z_\nu(b)=E_\nu$.
This proves the energy comparison.  The same proof with nonnegative
determinants applies for $m=0$; alternatively use
$\|H_m-H_0\|\le2Nm$ and finite-volume continuity.
Finally the commuting integer Gauss operators decompose
$\mathcal H$ into their joint eigenspaces.  The inequality for every
joint eigenvalue proves that the auxiliary ground energy equals $E_0$.
\end{proof}

This proof uses the full local Gauss projection, not just its global
charge component.  The constant gauge transformation acts trivially on
the links and imposes \eqref{lat-capacity} on the fermions.  The link
holonomies, including torus harmonic variables, remain integrated.

\begin{corollary}[Rest-energy strengthening in the physical sectors]
\label{lat-rest-energy}
Let $Q_\nu=\sum_x\nu_x$.  For the same model and every admissible
integer source,
\begin{equation}
 E_\nu\ge E_0+m|Q_\nu|.
 \label{lat-rest-energy-bound}
\end{equation}
For $m>0$ the corresponding thermal bound is
\begin{equation}
 Z_\nu(b)\le
 \binom{4N}{2N+|Q_\nu|}e^{-bm|Q_\nu|}Z_0(b).
 \label{lat-rest-energy-thermal}
\end{equation}
\end{corollary}
\begin{proof}
Choose a root site $x_*$ and decompose the Haar angles as
$\theta_x=\vartheta_x+\phi$, with $\vartheta_{x_*}=0$.
This is a measure-preserving change of torus coordinates.  The endpoint
$U^{e^{i\theta}}=U^{e^{i\vartheta}}$ and hence the bridge measure and
all scalar gauge weights are independent of $\phi$.  At a fixed
relative twist and fixed bridge, the remaining fermion factor is the
centered trace with global charge angle $\phi$ and one-particle history
$e^{i\vartheta}V_\omega$.

Lemma~\ref{lem:mass-gap-annulus} applies with $M(t)\ge mI$ and total
free time $b$.  It also applies to the continuous bridge history:
the same absolutely continuous identity
$\frac{d}{dt}\langle f,\beta f\rangle=-2\langle f,M(t)f\rangle$
gives its eigenvalue bounds directly.  Thus
Lemma~\ref{lem:mass-gap-fourier}, with one-particle dimension $4N$,
bounds the absolute value of its $Q_\nu$ Fourier coefficient by
$\binom{4N}{2N+|Q_\nu|}e^{-bm|Q_\nu|}$ times its nonnegative zeroth
coefficient.  Integrate first in $\phi$, use this bound, and remove the
remaining source phase $e^{-i\vartheta\cdot\nu}$ by absolute value.
The integral of the zeroth coefficient over relative twists and gauge
bridges is exactly the full local-Gauss vacuum partition function
$Z_0(b)$.  This proves \eqref{lat-rest-energy-thermal}.  The finite
binomial factor disappears as $b\to\infty$.  For $m=0$, the conclusion
is Theorem~\ref{lat-main}.
\end{proof}

The strengthening concerns the total electron excess forced by external
charges.  It is not a neutral-sector excitation gap or a photon gap.

\subsection{Nuclear self-energy and a uniform linear bound}

Let $D:\mathbb R^{\mathcal E}\to\mathbb R^X$ be the incidence
matrix, with $(DE)_x=\operatorname{div}E_x$, and let
$\Delta=DD^*$ be the nonnegative combinatorial lattice Laplacian.
Write $\mathcal G_n=\Delta^+$ for its inverse on zero-mean functions,
extended by zero on constants.  On $\mathcal H_\nu$, the exact identity
\begin{equation}
 H_E=\frac{g^2}{2a}
 \left[(\nu-q)^*\mathcal G_n(\nu-q)
       +\left\|E-D^*\mathcal G_n(\nu-q)\right\|^2\right]
 \label{lat-longitudinal-split}
\end{equation}
holds as a quadratic-form identity.  It follows from the orthogonal
decomposition of a real link vector into $\operatorname{ran}D^*$
and $\ker D$, since $DE=\nu-q$ has zero sum.  The commuting electric
and occupation operators permit the same computation in their joint
eigenbasis.  The second term is positive, but its allowed spectrum
depends on the integer-flux and charge constraints.  Equation
\eqref{lat-longitudinal-split} does not assert a tensor factorization
into independent transverse and longitudinal physical Hilbert spaces.

The source-source term in the first square is
\begin{equation}
 \frac{g^2}{2a}\nu^*\mathcal G_n\nu
 =\frac{g^2}{2a}\mathcal G_n(0)\sum_iZ_i^2
   +\frac{g^2}{a}\sum_{i<j}Z_iZ_j\mathcal G_n(R_i-R_j).
 \label{lat-source-self}
\end{equation}
Thus the conventional subtraction of individual nuclear self-energies
defines, explicitly,
\begin{equation}
 E_\nu^{\rm sub}=E_\nu-S_n(R,Z),\qquad
 S_n(R,Z)=\frac{g^2}{2a}\mathcal G_n(0)\sum_iZ_i^2.
 \label{lat-self-subtraction}
\end{equation}
The subtraction is prescribed per nuclear species and is independent
of the interacting state and of nuclear positions.  Its mild finite-box
dependence is fixed by the periodic Green function.  A prescribed
infinite-lattice self-energy could instead be used, but would be a
different, equally explicit normalization.

\begin{lemma}[Uniform lattice self-energy]
\label{lat-green-bound}
For every $n\ge2$ in three dimensions,
\begin{equation}
 0\le\mathcal G_n(0)
  =\frac1{n^3}\sum_{k\ne0}
       \frac1{4\sum_{j=1}^3\sin^2(\pi k_j/n)}\le\frac{13}{16}.
 \label{lat-green-diagonal}
\end{equation}
The sum runs over one reciprocal-lattice representative per residue.
\end{lemma}
\begin{proof}
Choose representatives with $|k_j|\le n/2$.  Since
$\sin(\pi|k_j|/n)\ge2|k_j|/n$, the denominator is at least
$16|k|^2/n^2$.  The shell $\max_j|k_j|=s$ contains at most
$(2s+1)^3-(2s-1)^3=24s^2+2\le26s^2$ points.  Therefore
\[
 \mathcal G_n(0)\le\frac1{16n}
      \sum_{0<\max|k_j|\le\lfloor n/2\rfloor}|k|^{-2}
 \le\frac{26\lfloor n/2\rfloor}{16n}\le\frac{13}{16}.
\]
For even $n$, the chosen representatives contain only a subset of some
boundary shells, which improves this upper bound.
\end{proof}

\begin{corollary}[Stability with subtracted individual source energies]
\label{lat-linear-stability}
For positive nuclei satisfying $\sum_i Z_i\le2n^3$, with any sites
$R_i$, and $1\le Z_i\le Z_{\max}$,
\begin{equation}
 E_\nu^{\rm sub}-E_0
 \ge m\sum_i Z_i-\frac{g^2}{2a}\mathcal G_n(0)\sum_iZ_i^2
 \ge m\sum_i Z_i-\frac{13g^2}{32a}Z_{\max}^2K,
 \label{lat-final-bound}
\end{equation}
where $13g^2/(32a)=13\pi\alpha/(8a)$ when $g^2=4\pi\alpha$.
The constant is uniform in box size, source positions, fermion pair
content, and electric or magnetic excitations.  It requires no
smallness of $\alpha$ and no restriction relating $a^{-1}$ to $m$.
\end{corollary}
\begin{proof}
Subtract the explicitly defined $S_n$ from
Corollary~\ref{lat-rest-energy} and use Lemma~\ref{lat-green-bound}.
\end{proof}

Two nuclei assigned to the same site remain two labeled sources in
\eqref{lat-self-subtraction}; their mutual energy
$g^2Z_iZ_j\mathcal G_n(0)/a$ is retained.  Hence the bound still
holds for lattice site collisions.  Replacing their individual
subtractions by the square of the total charge at that site removes
this mutual term and defines a different model.  With arbitrarily many
nuclei at a site, that alternative subtraction is not bounded by a
constant times $K$, so \eqref{lat-final-bound} does not follow from this
proof.  A lattice site is an exact source location for the lattice model,
but represents resolution of order $a$ when compared with continuum
point nuclei; no continuum collision singularity is being preserved.

\paragraph{A volume-independent source convention.}
One may instead prescribe the infinite-lattice diagonal
\[
\mathcal G_\infty(0)=
\int_{[-1/2,1/2]^3}
\frac{dt}{4\sum_j\sin^2(\pi t_j)}
\le\frac1{16}\int_{|t|\le\sqrt3/2}|t|^{-2}\,dt
=\frac{\pi\sqrt3}{8}.
\]
The integrable singularity in three dimensions makes this finite.
It is the zero-site Green value, equivalently the monotone
zero-mass resolvent integral of the infinite lattice Laplacian.
Subtracting the fixed per-species constants
\(S_\infty=\frac{g^2}{2a}\mathcal G_\infty(0)\sum_iZ_i^2\)
from the raw source-sector energy gives
\[
E_\nu-S_\infty-E_0
\ge mZ_{\rm tot}-\frac{\pi\sqrt3\,g^2}{16a}Z_{\max}^2K.
\]
This convention is independent of the box size as well as of the
state and positions. It is an alternative specified normalization,
not a claim that the finite-box diagonal equals its infinite-lattice value.

\subsection{Scope and limitations}

The substantive comparison is $E_\nu\ge E_0$ before self-energy
subtraction, not the elementary finite-volume lower bound.  After the
subtraction, it says precisely that nuclei together with their
specified individual electromagnetic self-energies cannot lower the
energy below the interacting vacuum.  Adding those self-energies to
the unsubtracted Hamiltonian would count them twice.

Unlike the sharp momentum regulator, this model is spatially local,
exactly gauge invariant, and has the extensive bound
\eqref{lat-extensive} as well as \eqref{lat-final-bound}.  These facts
alone do not establish a thermodynamic limit, atomic binding, a gap
above the vacuum in its neutral Gauss sector, or crystalline order.
The vacuum and open-source thermodynamic limits are proved by
additional arguments in Section~\ref{therm4-section}; the periodic-source
limit is distinguished there.
The charge-cost strengthening \eqref{lat-rest-energy-bound} does not
control neutral pair or photon excitations.

The electric energy here is the unshifted compact-rotor Laplacian.
A nonintegral background electric shift can make its kernel complex.
Restricting to an additional flux or winding sector can introduce a
further character into the gauge integral.  The proof has made neither
modification and does not compare separately fixed flux sectors.
In the presence of mobile charged matter, some proposed global electric
flux restrictions are not even reducing subspaces of the hopping
Hamiltonian.  A finite electric-field truncation also need not retain
the positive kernel \eqref{lat-positive-heat}.

Finally, positivity of the \emph{bare} mass in
\eqref{lat-wilson-mass} is a structural hypothesis.  Wilson fermions
have additive mass renormalization; fixed physical parameters need not
remain in this bare parameter range along a putative continuum
trajectory.  The source bound itself diverges as $a^{-1}$.
Neither a renormalized continuum QED construction nor a cutoff-uniform
stability constant follows.  The nonnegative determinant mechanism is
not a proof that compact \(U(1)\) is in a Coulomb phase at every coupling,
nor a proof of gapless transverse photons. It is
closely related to familiar positivity and static-source character
bounds in lattice gauge theory; no priority claim is made for that
mechanism or for this application.
